\chapter{\B{Preface}{Prefacio}}
\B{You own a Jamaican patty franchise in the first story. A patty is a baked pastry with a filling. Your team brings you three shops to choose from. You set the rules, compare the numbers and decide which shop to visit again before you sign a lease. In later chapters you will run a ferry desk, check a plan for loading a boat and lead a school club's AI demo.}
{En la primera historia eres dueño de una franquicia de patties en Jamaica. Un patty es una empanada horneada con relleno. Tu equipo te trae tres locales para elegir. Tú fijas las reglas, comparas las cifras y decides qué local volver a visitar antes de firmar. Más adelante estarás a cargo de una ventanilla de ferris, revisarás el plan de carga de un barco y dirigirás una demostración de IA en tu escuela.}

\B{A large language model, or LLM, learns patterns from examples of text. It can use those patterns to continue a message or draft an answer. This book shows you what happens between typing a question and reading the reply. It also shows what you can check when a reply affects a decision. The stories and their practice records are fictional. The public figures in Chapter~\ref{ch:capstone} come from the sources named there.}
{Un modelo grande de lenguaje, o LLM, aprende patrones a partir de ejemplos de texto. Puede usarlos para continuar un mensaje o preparar una respuesta. Este libro muestra qué pasa desde que escribes una pregunta hasta que lees la respuesta. También muestra qué puedes comprobar si una respuesta afecta una decisión. Las historias y sus registros de práctica son ficticios. Las cifras públicas del capítulo~\ref{ch:capstone} vienen de las fuentes que se indican allí.}

\B{Start with the story. Make a choice before you read the result. Follow the worked example with a pencil and a calculator, then explain the result in your own words. You can use the story and its decision even if algebra is new to you. Sections marked \emph{More maths} show how to write the same ideas with symbols. Return to them when you want the full calculation. The answers to the exercises are at the back.}
{Empieza con la historia. Toma una decisión antes de leer el resultado. Sigue el ejemplo resuelto con lápiz y calculadora. Luego explica el resultado con tus propias palabras. Puedes resolver la decisión de la historia aunque estés empezando con el álgebra. Las partes marcadas \emph{Más matemáticas} muestran cómo escribir las mismas ideas con símbolos. Vuelve a ellas cuando quieras hacer el cálculo completo. Las soluciones están al final del libro.}

\B{A ferry notice can give different times for different islands. A produce order needs a unit such as kilograms. A school message can switch languages in the same sentence. The book gives the meaning of short language samples beside the examples. A person may use another spelling or way of saying the same thing. Keep the message as written so people who know the language can check your answer.}
{Un aviso de ferri puede dar horas distintas según la isla. Un pedido de productos agrícolas necesita una unidad, como kilogramos. Un mensaje escolar puede cambiar de idioma en la misma oración. El libro explica el sentido de las frases breves junto a los ejemplos. Una persona puede usar otra grafía u otra manera de decir lo mismo. Guarda el mensaje tal como llegó para que quienes conocen el idioma puedan revisar tu respuesta.}

\B{Use the place index to find a case near you. It follows the United Nations list of Caribbean areas, called M49, and adds cases for named islands and places on the mainland. Some tasks ask you to check the island because the answer changes from one island to the next \citep{m49,caricom,acs}.}
{Usa el índice de lugares para buscar un caso cerca de ti. Sigue la lista de áreas del Caribe de las Naciones Unidas, llamada M49, y añade casos de islas y lugares del continente. Algunas tareas te piden comprobar la isla porque la respuesta cambia de una a otra \citep{m49,caricom,acs}.}

\B{Tong Xiao and Jingbo Zhu's textbook explains several of these methods in more depth. Look for it in the bibliography when you want to go further \citep{xiao}.}
{El libro de Tong Xiao y Jingbo Zhu explica varios de estos métodos con más detalle. Búscalo en la bibliografía si quieres profundizar \citep{xiao}.}

\chapter{\B{Notation}{Notación}}
\B{Use this page when you want to look up a symbol. You can start Chapter~\ref{ch:data} with the shop story and come back here when you need it.}
{Usa esta página cuando quieras buscar un símbolo. Puedes empezar el capítulo~\ref{ch:data} con la historia del negocio y volver aquí cuando lo necesites.}

\noindent\begin{tabularx}{\textwidth}{@{}>{$}l<{$}Y@{}}
\toprule
\text{\B{Symbol}{Símbolo}} & \B{Meaning in this book}{Significado en este libro}\\
\midrule
x,\;x_t & \B{Input sequence and token at position \(t\).}{Secuencia de entrada y token en la posición \(t\).}\\
x_{<t} & \B{Tokens before position \(t\).}{Tokens que preceden a la posición \(t\).}\\
\mathcal V,\;|\mathcal V| & \B{Tokenizer vocabulary and its size.}{Vocabulario del tokenizador y su tamaño.}\\
\theta,\;\theta_0 & \B{Model parameters and a reference set of parameters.}{Parámetros del modelo y conjunto de referencia.}\\
p_\theta(y\mid x) & \B{Probability the model assigns to \(y\) given \(x\).}{Probabilidad que el modelo asigna a \(y\) dado \(x\).}\\
z,\;\softmax(z) & \B{Scores before normalisation and the resulting probability vector.}{Puntuaciones previas a la normalización y vector de probabilidades resultante.}\\
Q,K,V & \B{Query, key and value matrices in attention. \(V\) differs from \(\mathcal V\).}{Matrices de consultas, claves y valores en atención. \(V\) se distingue de \(\mathcal V\).}\\
d_k & \B{Number of components in a query or key for one head.}{Número de componentes de una consulta o clave en una cabeza.}\\
M,\;M_{ij} & \B{Attention mask and entry controlling access from position \(i\) to \(j\).}{Máscara de atención y entrada que controla el acceso de la posición \(i\) a \(j\).}\\
\mathcal L,\;\eta & \B{Loss and learning rate.}{Pérdida y tasa de aprendizaje.}\\
\ln,\;\exp & \B{Natural logarithm and its inverse. \(\ln(1)=0\), and \(\exp(0)=1\).}{Logaritmo natural y su inversa. \(\ln(1)=0\) y \(\exp(0)=1\).}\\
\sum,\;\prod & \B{Add terms and multiply terms, respectively.}{Suma de términos y producto de términos, respectivamente.}\\
\argmax,\;\argmin & \B{Input yielding the largest or smallest score, respectively.}{Entrada que da la mayor o menor puntuación, respectivamente.}\\
\mathbf 1[\cdot] & \B{Indicator. It is 1 when its condition holds and 0 otherwise.}{Indicador. Vale 1 si se cumple la condición y 0 en caso contrario.}\\
\text{MiB},\;\text{GiB} & \B{\(2^{20}\) and \(2^{30}\) bytes.}{\(2^{20}\) y \(2^{30}\) bytes.}\\
\bottomrule
\end{tabularx}

\medskip
\noindent\B{Lowercase letters usually denote scalars or sequences. A vector is an ordered list of numbers. A matrix arranges numbers in rows and columns. Superscript \({}^\top\) transposes a matrix, exchanging rows and columns. The context of an equation identifies when a symbol is reused for a different quantity.}
{Las letras minúsculas suelen indicar escalares o secuencias. Un vector es una lista ordenada de números. Una matriz organiza números en filas y columnas. El superíndice \({}^\top\) transpone una matriz, intercambiando filas y columnas. El contexto de cada ecuación indica si un símbolo se reutiliza para otra magnitud.}
